\chapter{Introduction} \label{chap:introduction}

Dans ce modèle, les styles sont paramétrés afin de respecter les normes de mise en page d’un mémoire ou d'une thèse présentées dans le \href{https://www.usherbrooke.ca/geomatique/fileadmin/sites/geomatique/espace-etudiant/guide-redaction2004_vf.pdf}{guide de rédaction de Jean-Marie-Dubois}. \textcolor{eclipsepurple}{\textbf{La table des matières et les titres des chapitres et des sections présentés dans ce modèle sont fournis à titre indicatif seulement. Ils doivent être revus et adaptés afin de refléter fidèlement les objectifs, la méthodologie et les résultats propres à votre recherche.}}


\section{Mise en contexte} \label{sec:mise_en_contexte}
Voici un paragraphe.

Voici un second paragraphe.

\lipsum[1-4] % génère 4 paragraphes (à supprimer)

\section{Problématique} \label{sec:problematique}

\lipsum[1-2] % génère 2 paragraphes (à supprimer)

\section{Objectifs} \label{sec:objectifs}

Décrire l’objectif principal et les objectifs spécifiques.

Bonnes pratiques pour nommer les \textit{labels} dans un document \LaTeX :

\begin{itemize}
    \item \textbf{chap:} pour les \textit{chapitres} : \verb|\label{chap:introduction}|.
    \item \textbf{sections :}
    \begin{itemize}
        \item \textbf{sec:} pour les \textit{sections de niveau 1} : \verb|\label{sec:mise_en_contexte}|.
        \item \textbf{subsec:} pour les \textit{sous-sections de niveau 2} : \verb|\label{subsec:sous_section}|.
        \item \textbf{subsubsec:} pour les \textit{sous-sections de niveau 3} : 
        \lstinline|\label{subsubsec:sous_sous_section}|.
    \end{itemize}
    \item \textbf{fig:} pour les \textit{figures} : \verb|\label{fig:territoire_etude}|.
    \item \textbf{tab:} pour les \textit{tableaux} : \verb|\label{tab:statistiques_descrip}|.
    \item \textbf{eq:} pour les \textit{équations} : \verb|\label{eq:rmse}|.
    \item \textbf{algo:} pour les \textit{algorithmes} : \verb|\label{algo:dijkstra}|.
\end{itemize}

Ces conventions permettent d'assurer une bonne organisation et lisibilité du code source, et d'éviter les conflits entre les \textit{labels}. Elles permettent aussi de faire référence à un élément dans le texte. Par exemple...

Comme discuté dans la section~\ref{sec:mise_en_contexte}, ...
